% Notes of the Montreal seminar, 1970 (cote 129, typed part): archive cote 129, site page 51 (PDF page 52).
% typescript. Draft compared: none.
% AI-transcribed by Claude Opus 5.5 (claude-opus-5-5, adaptive thinking, effort high), Message Batches API, batch mtl129_full (msgbatch_01AGnDKHwy6ATCM7ehUa6263), 2026-11-13; not yet checked by eye.
\scanpage{129/51}
\mspage{\hand{II 23}}

$V^{N+1}(u(x)) = 0$ car $V^{N+1}(x) = 0$ ; l'élément $u(x)$ provient par la
limite inductive d'un élément $(y_1,y_2,\ldots,y_r) \in \mathbb{W}_r(S)$ et on a
$V^{N+1}(y_1,\ldots,y_r) = 0$ car $\mathbb{W}_r$ se plonge dans $\underset{\rightarrow}{\mathbb{W}}$ ; on a donc
$y_1 = y_2 = \ldots y_{r-N-1} = 0$ , ce qui implique que $(y_1,y_2,\ldots,y_r)$ pro-%
vient d'un élément unique de $\mathbb{W}_{N+1}(S)$ par le morphisme de transition
$T^{r-N-1}$ .

On en déduit que tout homomorphisme $u \in \mathrm{Hom}_{\mathcal{D}}(M, \mathbb{W}(S))$ se
factorise de manière unique par $W_{N+1}(S) \hookrightarrow \underset{\rightarrow}{\mathbb{W}}(S)$ .

On est donc ramené au problème de construire un isomor-%
phisme fonctoriel en $S$
\[ \phi_S : \mathrm{Hom}_{\mathcal{D}}(M, \mathbb{W}_{N+1}(S)) \rightarrow \mathrm{Hom}_k(A, \Gamma(\underline{\mathcal{O}}_S)) . \]

Pour tout $u \in \mathrm{Hom}_{\mathcal{D}}(M, \mathbb{W}_{N+1}(S))$ , $\phi_S(u) : A \rightarrow \Gamma(\mathcal{O}_S)$ est
le morphisme qui associe à $T_x$ la dernière coordonnée de
\margin{$\times$} $u(x) \in \mathbb{W}_{N+1}(S) = \Gamma(\mathcal{O}_S)^{N+1}$ . Il est clair que $\hand{\phi}_S(u)$ est bien défini
sur $A$ et est un morphisme d'algèbres (à cause des relations vérifiées
dans $A$). On définit une application inverse en associant à $\overline{u} \in \mathrm{Hom}_k(A, \Gamma(\mathcal{O}_S))$
l'homomorphisme de $\mathcal{D}$-modules $u : M \rightarrow W_{N+1}(S)$ qui à $x$ fait corres-%
pondre $u(x) = (\overline{u}(T_{V^N(x)})\ldots, \overline{u}(T_x))$ .

6.3. \uline{Le schéma} $\mathbb{E}(M)$ \uline{est un p-groupe fini unipotent}

Par définition même de $\mathbb{E}(M)$ (6.1), $\mathbb{E}(M)$ est un schéma
en groupes, dépendant additivement de $M$ . Comme $M$ est annulé par une
puissance de $p$ (car $V$ est nilpotent), il en est de même de $\mathbb{E}(M)$ .
Montrons que $\mathbb{E}(M)$ est fini.

\mspage{\hand{51}}
